% Pursuing Stacks (cotes 134-2 to 134-8): archive cote 134-5, site page 80 (PDF page 81), typescript page 267.
% typescript. Draft compared: ps4.tex:4077-4157.
% AI-transcribed by Claude Opus 5.5 (claude-opus-5-5, adaptive thinking, effort high), Message Batches API, batch ps2 (msgbatch_01NVwr1Yuo8qJf5rykeGBXPK), 2026-11-12; not yet checked by eye.
\scanpage{134-5/80}
\mspage{267}

26.6.

82. Yesterday I have been busy mainly with the readjustment of \xout{\unsure{general}} \inserted{the overall}
perspective on the main notions developped so far, which has sprung from
the new fact stated at the beginning of yesterday's notes: namely that an
arbitrary contractibility structure $M_c$ (on a category M stable under finite
products) can be recovered in terms of the associated $\underline{W}$-asphericity struc-%
ture, where $\underline{W}$ is any basic localizer satisfying (besides the conditions
Loc 1) to Loc 3) of p. 213-214) the extra assumption Loc 4) of p.260. It
seems about time now to enter into some \xout{technical} a little more technical
specifications along the same lines - and to start with, give a proof
of this ``new fact'' ! Let's state it again in full:

\uline{Theorem} 1 . Let M be a $\underline{U}$-category stable under finite products, $M_c$ \inserted{$\subset$ Ob M} a con-%
tractibility structure on M, admitting a small full subcategory C which
generates the structure. Let moreover $\underline{W}$ be any basic localizer satisfying
condition Loc 4) (compatibility with the $\underline{\pi}_0$-functor (Cat)$\longrightarrow$(Sets)), and
let $M_{\underline{W}} \subset$ Ob M the $\underline{W}$-asphericity structure generated by $M_c$ , $W = \underline{W}_{M_c}$ the
corresponding set of ``$\underline{W}$-equivalences'' or ``weak equivalences'' in \unsure{Fl}(M). Con
-sider the homotopy structure $h_W$ associated to W, i.e. (cf. section 54)
the homotopy \xout{interval} structure associated to the homotopy interval struc-%
ture \inserted{J} generated by the set \inserted{$J_0$} of all intervals
\[ \underline{I} = (I, \hand{\delta}_0, \hand{\delta}_1) \]
in M such that
\[ I \hand{\in} M_{\underline{W}} \qquad . \]
Then $h_W$ \xout{\illegible{1}} is the homotopy structure on M associated to the contractibi-%
lity structure $M_c$ , and hence $M_c$ can be described in terms of $M_{\underline{W}}$ (or of
$\underline{W}_{M_c} = W$ ) as the set of objects which are contractible for $h_W$ , i.e. such
that the map $x \longrightarrow e_M$ is an $h_W$-homotopism.

\uline{Proof}. Let $M'_c$ be the set of $h_W$-contractible objects of M, clearly we have
\[ M_c \subset M'_c \subset M_{\underline{W}} \overset{\mathrm{def}}{=} M_a . \]
The theorem amounts to saying that $M_c$ generates the homotopy interval
structure J (by which we mean that the set of intervals of M made up with
objects of $M_c$ generates the structure J). Indeed, because of $M_c \subset M'_c$
this will imply that J is associated to a contractibility structure, \xout{namely}
\xout{ly} namely to $M'_c$ . But for an object x of M to be in $M'_c$ i.e. to be contrac
tible for the structure J, amounts to be contractible for $M_c$ , and hence
by saturation of $M_c$ , to be in $M_c$ , hence $M'_c = M_c$ - which yields what we

\margin{80}
